\newif\ifglobal

%%%%%%%%%%%%%%%%%%%%%%%
%%% Compile to view %%%
%%%%%%%%%%%%%%%%%%%%%%%

%%%%%%%%%%%%%%%%%%%%%%%%%%%%%%%%%%%%%%%%%%%%%%%%%%%
%%% Readme:                                     %%%
%%% https://www.overleaf.com/read/bwdjvfjksvjy/ %%%
%%%%%%%%%%%%%%%%%%%%%%%%%%%%%%%%%%%%%%%%%%%%%%%%%%%

% >>> M?: marks a module setting number or important notice

% >>> M4: Set {./relative-path} to external files for this module <<<
\def \modulefiles {./07-BOOTCTRL}
% To add an external file, use:
% (Replace '---' with actual filename)
% '\input{\modulefiles/tables/---.tex}' for tables
% '\includegraphics{\modulefiles/figures/---}' for figures
% '\subfile{\modulefiles/---}' for register files

% >>> M1: This module is in global project? <<<
% 'YES' - uncomment; 'NO' - comment out
\globaltrue

\ifglobal
    \documentclass[main\_\_EN.tex]{subfiles}
   
\else
    \documentclass[openany, 10pt]{book}

    % >>> M2: In ./00-shared/preamble-trm-module, also find
    %         settings for visibility of todo notes, watermark <<<
    \usepackage{preamble-trm-module}

    % >>> M3: Temporary labels in English,
    %         you can add your own labels there <<<
    \myexternaldocument{./00-shared/temp-labels-en}

\fi %\ifglobal


\setcounter{secnumdepth}{4}


% >>> M5: If this module needs any updates in
% './00-shared/preamble-trm-module.sty'
% (include additional package, etc.), contact your module PM
% or create the file './00-shared/changelog.tex' make the
% necessary updates yourself and document them in this file <<<

% Make text left-aligned and allow latex to break pages naturally, comment out for full justification
\raggedright
\raggedbottom


\begin{document}


% Sets language into which pre-defined bits of text are translated
\selectlanguage{english}
\setenglishformatting

% Do not delete this block
\ifglobal
% Add the following front matter if
% this module is outside of global project
\else
    \InputIfFileExists{readme}{}{}
    \listoftodos
    {\let\clearpage\relax \listofdonetodos}
    \tableofcontents
    %\listoftables
    %\listoffigures
\fi



%%%%%%%%%%%%%%%%%%%%%%%%%
%%% TRM Module Begins %%%
%%%%%%%%%%%%%%%%%%%%%%%%%

\hypertarget{bootctrl}{}
\chapter{Chip Boot Control (BOOTCTRL)}
\label{mod:bootctrl}

\section{Overview}

ESP32-S2 has three strapping pins:
\begin{itemize}
    \item GPIO0
    \item GPIO45
    \item GPIO46
\end{itemize}

Software can read the values of the three strapping pins from register \hyperref[fielddesc:GPIOSTRAPPING]{GPIO\_STRAPPING}.
During power-on-reset, RTC watchdog reset, brownout reset, analog super watchdog reset, and crystal clock glitch detection reset (see Chapter \ref{mod:resclk} \textit{\nameref{mod:resclk}}), hardware samples and stores the voltage level of strapping pins as strapping bit of “0” or “1” in latches, and hold these bits until the chip is powered down or shut down.

By default, GPIO0, GPIO45 and GPIO46 are connected to internal pull-up/pull-down during chip reset. Consequently, if the three GPIOs are unconnected or their connected external circuits are high-impedance, the internal weak pull-up/pull-down will determine the default input level of the strapping pins (see Table \ref{table:pull_up_down}).

\begin{longtable}[c]{|p{2.5cm}|p{2.5cm}|p{2.5cm}|p{2.5cm}|}
\caption{Default Configuration of Strapping Pins}
\label{table:pull_up_down}
    \\\hline
    \rowcolor{lightgray}
  Pin & GPIO0 & GPIO45 & GPIO46 \\ \hline
  Default & Pull-up & Pull-down & Pull-down \\ \hline
\end{longtable}

To change the default configuration of strapping pins, users can apply external pull-down/pull-up resistors, or use host MCU GPIOs to control the voltage level of these pins when powering on ESP32-S2. After the reset is released, the strapping pins work as normal-function pins.

\begin{tiplisting}
The following section provides description of the chip functions and the pattern of the strapping pins values to invoke each function. Only documented patterns should be used. If some pattern is not documented, it may trigger unexpected behavior.
\end{tiplisting} 

\section{Boot Mode}

GPIO0 and GPIO46 control the boot mode after reset.

\begin{longtable}[c]{|p{2.5cm}|p{2.5cm}|p{2.5cm}|}
\caption{Boot Mode}
\label{table:boot_mode}
    \\\hline
    \rowcolor{lightgray}
  Pin & SPI Boot & Download Boot \\ \hline
  GPIO0 & 1 & 0 \\ \hline
  GPIO46 & x & 0 \\ \hline
\end{longtable}

Table \ref{table:boot_mode} shows the strapping pin values and the associated boot modes. "x" means that this value is ignored. Currently only the two boot modes shown are supported. The strapping combination of GPIO0 = 0 and GPIO46 = 1 is not supported and will trigger unexpected behavior.

In SPI boot mode, the CPU boots the system by reading the program stored in SPI flash.

In download boot mode, users can download code to SRAM or flash using UART0, UART1, QPI or USB interface. It is also possible to load a program into SRAM and execute it in this mode.

The following eFuses control boot mode behavior:

\begin{itemize}
    \item \hyperref[fielddesc:EFUSEDISFORCEDOWNLOAD]{EFUSE\_DIS\_FORCE\_DOWNLOAD} 
    \begin{itemize}
        \item If this eFuse is 0 (default), software can switch the chip from SPI boot mode to download boot mode by setting register \hyperref[fielddesc:RTCCNTLFORCEDOWNLOADBOOT]{RTC\_CNTL\_FORCE\_DOWNLOAD\_BOOT} and triggering a CPU reset. In this case, hardware overwrites \hyperref[fielddesc:GPIOSTRAPPING]{GPIO\_STRAPPING}[3:2] from ``1x" to ``00".
        \item If this eFuse is 1, this register is disabled. \hyperref[fielddesc:GPIOSTRAPPING]{GPIO\_STRAPPING} can not be overwritten. 
    \end{itemize}    
    \item \hyperref[fielddesc:EFUSEDISDOWNLOADMODE]{EFUSE\_DIS\_DOWNLOAD\_MODE}
 
    If this eFuse is 1, download boot mode is disabled. \hyperref[fielddesc:GPIOSTRAPPING]{GPIO\_STRAPPING} will not be overwritten by \hyperref[fielddesc:RTCCNTLFORCEDOWNLOADBOOT]{RTC\_CNTL\_FORCE\_DOWNLOAD\_BOOT}.
    \item \hyperref[fielddesc:EFUSEENABLESECURITYDOWNLOAD]{EFUSE\_ENABLE\_SECURITY\_DOWNLOAD}


If this eFuse is 1, download boot mode only allows reading, writing and erasing plaintext flash and does not support any SRAM or register operations. This eFuse is ignored if download boot mode is disabled.
\end{itemize}

\section{ROM Messages Printing Control}

During the boot process, the messages by the ROM code can be printed to:

\begin{itemize}
    \item \textbf{(Default) UART0}
    \item UART1
\end{itemize}

\hyperref[fielddesc:EFUSEUARTPRINTCHANNEL]{EFUSE\_UART\_PRINT\_CHANNEL} controls if the ROM messages will be printed to UART0 or UART1.
\begin{itemize}
    \item 0: UART0
    \item 1: UART1
\end{itemize}

\hyperref[fielddesc:EFUSEUARTPRINTCONTROL]{EFUSE\_UART\_PRINT\_CONTROL} and GPIO46 control ROM messages printing to \textbf{UART}.


\begin{table}[H]
    \centering
    \caption{UART ROM Message Printing Control}
    \label{table:ROM_CODE_CONTROL}

    \begin{threeparttable}

    \begin{tabular}{ | l | c | C{1.3cm} |}
    \hline

    \rowcolor{lightgray}
    \textbf{UART ROM Code Printing} & \textbf{EFUSE\_UART\_PRINT\_CONTROL} &   \textbf{GPIO46}  \\\hline

    \multirow{3}{*}{\textbf{Enabled}} & \textbf{0} & Ignored \\ \cline{2-3}
     & 1 & 0 \\ \cline{2-3}
     & 2 & 1 \\ \hline
    \multirow{3}{*}{Disabled} & 1 & 1 \\ \cline{2-3}
     & 2 & 0 \\ \cline{2-3}
     & 3 & Ignored \\ \hline

    \end{tabular}

        \begin{tablenotes}
            \item[1] \textbf{Bold} marks the default value and configuration.
        \end{tablenotes}

    \end{threeparttable}
\end{table}



\section{VDD\_SPI Voltage}

GPIO45 is used to select the VDD\_SPI power supply voltage at reset:

\begin{itemize}
    \item GPIO45 = 0, VDD\_SPI pin is powered directly from VDD3P3\_RTC\_IO via resistor R$_{SPI}$. Typically this voltage is 3.3 V.
    
    \item GPIO45 = 1, VDD\_SPI pin is powered from internal 1.8 V LDO.
\end{itemize}

This functionality can be overridden by setting \hyperref[fielddesc:EFUSEVDDSPIFORCE]{EFUSE\_VDD\_SPI\_FORCE} to 1, in which case the \hyperref[fielddesc:EFUSEVDDSPITIEH]{EFUSE\_VDD\_SPI\_TIEH} value determines the VDD\_SPI voltage:

\begin{itemize}
    \item \hyperref[fielddesc:EFUSEVDDSPIFORCE]{EFUSE\_VDD\_SPI\_FORCE} = 1 and \hyperref[fielddesc:EFUSEVDDSPITIEH]{EFUSE\_VDD\_SPI\_TIEH} = 0, VDD\_SPI connects to 1.8 V LDO.
    \item \hyperref[fielddesc:EFUSEVDDSPIFORCE]{EFUSE\_VDD\_SPI\_FORCE} = 1 and \hyperref[fielddesc:EFUSEVDDSPITIEH]{EFUSE\_VDD\_SPI\_TIEH} = 1, VDD\_SPI connects to VDD3P3\_RTC\_IO.
\end{itemize}

\end{document}
